\chapter*{Referee Verification Checklist}
\addcontentsline{toc}{chapter}{Verification Checklist}

To facilitate peer review, we provide a checklist of critical components, formal dependencies, and verification status.

\section*{1. Critical Lemma Dependencies}

\begin{itemize}
    \item \textbf{Lemma 2.1 (Reflection Positivity):} 
    \begin{itemize}
        \item Status: \textbf{Proved}.
        \item Dependency: Standard Group Integrals.
        \item Location: Chapter 2.
    \end{itemize}
    
    \item \textbf{Lemma 3.4 (Polymer Activity Bound):}
    \begin{itemize}
        \item Status: \textbf{Proved (Strong Coupling)}.
        \item Constants: Explicit $C(\beta)$ provided.
        \item Location: Chapter 3.
    \end{itemize}

    \item \textbf{Lemma 4.2 (Uniform Log-Sobolev Constant):}
    \begin{itemize}
        \item Status: \textbf{Conditional}.
        \item Hypothesis: Assumes bounded local specific heat.
        \item Location: Chapter 4.
    \end{itemize}
\end{itemize}

\section*{2. Circularity Checks}

\begin{itemize}
    \item \textbf{Check:} Does the continuum limit derivation assume the mass gap a priori?
    \item \textbf{Verification:} No. The existence of the limit is derived from Mosco convergence of the forms. The spectral gap is a property of the limiting form, deduced from uniform lower bounds on the lattice.
    
    \item \textbf{Check:} Is the scale setting $\beta(a)$ input or output?
    \item \textbf{Verification:} The scaling trajectory is defined by fixing the string tension $\sigma$. The observable $\sigma a^2$ is monotonic in $\beta$ (Strong Coupling) and assumed monotonic globally (Hypothesis $H_{Mono}$).
\end{itemize}

\section*{3. Explicit Constants}

\begin{table}[h]
\centering
\begin{tabular}{|l|l|l|}
\hline
\textbf{Constant} & \textbf{Value/Bound} & \textbf{Region Validity} \\ \hline
$\beta_0$ (Radius of Conv.) & $\beta < 0.2$ (approx) & Strong Coupling \\ \hline
$m_0$ (Mass Gap) & $\sim -4 \ln \beta$ & Strong Coupling \\ \hline
$c_{LSI}$ (LSI const) & $ > 0$ & Conditional \\ \hline
\end{tabular}
\end{table}
